\begin{theorem}\label{thm:tilts}
Fix $0<\rho<1$, $1<s<s_c:=(1+\rho)/(1-\rho)$ and $B>1$.
There are strictly positive mean-one densities $r_L$, for every $L\ge1$,
with $\|r_L\|_s\le B$ and $r_L\ge1/B$, and constants
$\kappa>1$ and $0<\omega<\kappa$, independent of $L$, such that
\[
 0<\gap(T_{r_L})\le
 2(\omega/\kappa)^L+\varepsilon^{-1}\kappa^{-L},
 \qquad \varepsilon=(B-1)/2.
\]
Thus a uniform positive lower density does not prevent exponential gap
collapse throughout the strict subcritical moment regime.
\end{theorem}
\begin{proof}
Relabel the cube as $\{-1,1\}^L$. For a fixed tilt $0<\delta<1$, put
\[
 n=\left(\frac{(1+\delta)^s+(1-\delta)^s}{2}\right)^{1/s},
 \qquad w_\pm(x)=\frac{1\pm\delta x}{n},\qquad m=n^{-1}.
\]
Strict convexity gives $n>1$, while $\|w_\pm\|_s=1$ and
$\mathbb E w_\pm=m$. Define $R_{\pm,L}(x)=\prod_i w_\pm(x_i)$ and
\[
 Z_L=1+2\varepsilon m^L,\qquad
 r_L=\frac{1+\varepsilon(R_{+,L}+R_{-,L})}{Z_L}.
\]
Product integration gives $\|R_{\pm,L}\|_s=1$ and
$\mathbb E R_{\pm,L}=m^L$. Hence $r_L>0$, $\mathbb E r_L=1$,
$\|r_L\|_s\le(1+2\varepsilon)/Z_L\le B$, and
$r_L\ge1/Z_L\ge1/B$.

Write $Q=P_{\sqrt\rho}$. Since $Q^2=P_\rho$, the matrices
$T_{r_L}=AA^*$ and $S_L=A^*A$, with $A=M_{\sqrt{r_L}}Q$,
have identical spectra, including multiplicities. Here $M_f$ denotes
multiplication by $f$. Direct expansion gives
\[
 Z_L S_L=P_\rho+\varepsilon(X_L+Y_L),\qquad
 X_L=H_+^{\otimes L},\quad Y_L=H_-^{\otimes L},
\]
where the one-site matrices, in the orthonormal Walsh basis $(1,x)$, are
\[
 H_\pm=\frac1n
 \begin{pmatrix}1&\pm\delta\sqrt\rho\\
                 \pm\delta\sqrt\rho&\rho\end{pmatrix}.
\]
They are positive definite and conjugate by $\diag(1,-1)$. Their common
simple top eigenvalue is
\[
 \kappa(\delta)=
 \frac{1+\rho+\sqrt{(1-\rho)^2+4\rho\delta^2}}{2n(\delta)}.
\]
At zero, $\kappa(0)=1$, $(\log\kappa)'(0)=0$, and differentiation yields
\[
 (\log\kappa)''(0)=\frac{2\rho}{1-\rho}-(s-1)>0.
\]
All functions are smooth near zero. Taylor's theorem therefore gives
$\kappa>1$ for every sufficiently small positive $\delta$. Fix one such
$\delta<1$. This choice depends on $\rho,s$ and is independent of $L$;
no uniform exponential rate as $s\uparrow s_c$ is asserted.

Set $\omega=\|H_+^{1/2}H_-^{1/2}\|$. Positive definiteness and
submultiplicativity give $0<\omega\le\kappa$. If equality held, norm
attainment at a unit vector $v$ would force equality in both steps of
\[
 \|H_+^{1/2}H_-^{1/2}v\|
 \le\sqrt\kappa\,\|H_-^{1/2}v\|\le\kappa.
\]
The second equality makes $v$ a top eigenvector of $H_-$; the first
then makes it a top eigenvector of $H_+$. But $H_+-H_-$ has zero
diagonal and nonzero equal off-diagonal entries, so is invertible.
It cannot annihilate their common top eigenvector. Thus $\omega<\kappa$.

For $N=2^L\ge2$, use the $N$-by-$2N$ matrix
$C=[\sqrt{X_L},\sqrt{Y_L}]$. Its two Gram matrices satisfy
\[
 CC^*=X_L+Y_L,\qquad
 C^*C=\begin{pmatrix}
 X_L&\sqrt{X_L}\sqrt{Y_L}\\
 \sqrt{Y_L}\sqrt{X_L}&Y_L
 \end{pmatrix}.
\]
The block diagonal part has its two largest eigenvalues equal to
$\kappa^L$. Tensorization of positive square roots and of operator
norms gives $\|\sqrt{X_L}\sqrt{Y_L}\|=\omega^L$, which is also the
norm of the self-adjoint off-diagonal perturbation. The variational
eigenvalue bound consequently gives
\[
 \lambda_1(X_L+Y_L)\le\kappa^L+\omega^L,\qquad
 \lambda_2(X_L+Y_L)\ge\kappa^L-\omega^L>0.
\]
Indeed $C^*C$ and $CC^*$ have the same $N$ positive eigenvalues,
with $N$ additional zeros only in $C^*C$, so both leading bounds
transfer even when $L=1$. Also $X_L+Y_L\ge X_L$ gives
$\lambda_1(X_L+Y_L)\ge\kappa^L$.

Finally $P_\rho$ is positive semidefinite of norm one. For
$F_L=Z_L S_L$, the variational principle and norm bound imply
\[
 \lambda_1(F_L)\le\varepsilon(\kappa^L+\omega^L)+1,
 \qquad \lambda_2(F_L)\ge\varepsilon(\kappa^L-\omega^L),
 \qquad \lambda_1(F_L)\ge\varepsilon\kappa^L.
\]
Subtracting and dividing gives the claimed upper bound. The scalar
$Z_L$ cancels from the ratio, and $S_L$ and the actual $T_{r_L}$
have the same spectrum. Its entrywise positivity gives strict Perron
separation at every finite $L$. Both terms in the upper bound decay
exponentially, although the bound can exceed one at small $L$.
\end{proof}

Theorems~\ref{thm:moment} and~\ref{thm:tilts} separate the two strict
regimes at the same value $s_c$. The critical case $s=s_c$ remains open
here. At $B=1$, Jensen forces $r=1$ and the actual gap is $1-\rho$;
the sufficient lower bound is not claimed optimal. The tilted family
does not claim product marginals for its weighting or stationary law.
It isolates the moment obstruction for the same transfer operator,
whereas Theorem~\ref{thm:main} isolates the marginal-information
obstruction. The earlier two-point-spike construction in the retained
analytical sources covers the narrower interval
$1<s<\log2/\log(2/(1+\rho))$; it remains valid but is superseded here
by the product-tilt construction.

